% Generated by jobs/build_results_table.py -- do not hand-edit.
% Table 1: comparison. Frozen TRELLIS.2 is also row 1 of the ablation table and MUST
% carry identical numbers there.
\newcommand{\cmark}{\ding{51}}   % requires \usepackage{pifont}
\newcommand{\xmark}{\ding{55}}
\begin{table}[t]
\centering
\small
\setlength{\tabcolsep}{4.5pt}
\begin{tabular}{lcccccc}
\toprule
& & \multicolumn{2}{c}{Reconstruction} & \multicolumn{3}{c}{Temporal (texel-space)} \\
\cmidrule(lr){3-4}\cmidrule(lr){5-7}
Method & Temporal & PSNR\,$\uparrow$ & SSIM\,$\uparrow$ & Flicker\,$\downarrow$ & Accel.\,$\downarrow$ & Drift \\
\midrule
Frozen TRELLIS.2 & \xmark & 11.50 & 0.3363 & 0.02074 & 0.03102 & 0.3621 \\
\midrule
Ours & \xmark & 24.50 & 0.7536 & 0.00839 & 0.01148 & 0.2213 \\
Ours & \cmark & \textbf{24.55} & \textbf{0.7554} & \textbf{0.00707} & \textbf{0.00771} & 0.2221 \\
\bottomrule
\end{tabular}
\caption{\textbf{Comparison against the frozen generator.} Averages over
24 objects. The frozen generator is not asked to produce a temporally coherent
texture, so the fidelity gap is large and uninformative on its own; the temporal columns
carry the claim. Ours is better on 24/24 objects for flicker and
24/24 for acceleration (two-sided sign test, $p < 10^{-6}$ for both),
reducing them 65.9\% and 75.1\%. Enabling the
temporal component accounts for 15.7\% of the flicker reduction and
32.8\% of the acceleration reduction on its own, at
+0.05\,dB. Drift is a guard rather than a score and is
never marked best: it moves +0.4\% when the temporal component is
enabled, confirming the gain is not damping the texture toward a still image --- the
degenerate solution that minimises both other temporal metrics. Component-wise attribution
is deferred to Table~\ref{tab:ablation}.}
\label{tab:comparison}
\end{table}
